\documentclass[11pt,a4paper]{book}

\usepackage[margin=2.6cm]{geometry}
\usepackage{amsmath,amssymb,amsthm}
\usepackage{stmaryrd}
\usepackage{enumitem}
\usepackage{booktabs}
\usepackage{tikz}
\usepackage{tikz-cd}
\usepackage[colorlinks=true,linkcolor=blue,citecolor=blue,urlcolor=blue]{hyperref}

%% ---- Notation (kept consistent with the rest of the repo) ----
\newcommand{\Set}{\mathbf{Set}}
\newcommand{\Cat}{\mathbf{Cat}}
\newcommand{\Cof}{\mathbf{Cof}}
\newcommand{\Cont}{\mathbf{Cont}}
\newcommand{\DCont}{\mathbf{DCont}}
\newcommand{\Poly}{\mathbf{Poly}}
\newcommand{\cont}[2]{#1 \triangleleft #2}
\newcommand{\Ext}[1]{\llbracket #1 \rrbracket}
\newcommand{\id}{\mathrm{id}}
\newcommand{\cod}{\mathrm{cod}}
\newcommand{\dom}{\mathrm{dom}}
\newcommand{\ob}{\mathrm{ob}\,}
\newcommand{\sh}{\sharp}
\newcommand{\co}{\mathrel{\text{\reflectbox{$\rightharpoonup$}}}}
\DeclareMathOperator{\Hom}{Hom}
\DeclareMathOperator{\End}{End}
\newcommand{\Cathash}{\mathbb{C}\mathrm{at}^\sharp}

%% ---- Flagging environments (Neil's convention) ----
\theoremstyle{plain}
\newtheorem{Theorem}{Theorem}[chapter]
\newtheorem{Proposition}[Theorem]{Proposition}
\newtheorem{Lemma}[Theorem]{Lemma}
\newtheorem{Corollary}[Theorem]{Corollary}
\newtheorem{Conjecture}[Theorem]{Conjecture}
\theoremstyle{definition}
\newtheorem{Definition}[Theorem]{Definition}
\newtheorem{Example}[Theorem]{Example}
\theoremstyle{remark}
\newtheorem{Remark}[Theorem]{Remark}

%% A small macro for the mandatory provenance tag, so it is visually uniform.
\newcommand{\prov}[1]{\hfill\textsf{\footnotesize[#1]}}

\title{\textbf{The Category of Containers}\\[0.4em]
  \large a living reference}
\author{MacBeth}
\date{Draft, June 2026, v0.1}

\begin{document}
\maketitle

\frontmatter
\tableofcontents

\chapter{Preface}
\markboth{PREFACE}{PREFACE}

\section*{Purpose of this document}

This is a single, evolving reference centred on \emph{the category of
containers}. It is meant to grow: chapters are added and revised in place,
and the whole document compiles to one PDF that can be read end to end or
consulted by section. The organising thesis is the equivalence chain
\[
  \textbf{Containers} \ \simeq\ \textbf{Directed Containers}
  \ \simeq\ \textbf{Polynomial Comonoids} \ \simeq\ \textbf{Small Categories},
\]
which we take as the foundation for \emph{verified composition}: once a
system (a supply chain, an agent hierarchy, a state space, a tax computation)
is presented as a container, its composition is governed by the structure of
$\Cont$ and of the comonoids inside it, and correctness of composition becomes
a statement that can be proved --- and, increasingly, machine-checked. Object
level equivalence is settled groundwork; the live mathematics lives at the
\emph{morphism} level, and that is where this book spends most of its effort.

\section*{The flagging convention}

Mathematical honesty is the first principle of this document. Every
\textbf{Definition} and every result is flagged, and every
\textbf{Theorem}/\textbf{Proposition}/\textbf{Lemma}/\textbf{Corollary}
statement \emph{ends with a bracketed provenance tag}, one of:
\begin{itemize}[nosep]
  \item \textsf{[Cited: \textit{author year, ref}]} --- a known result from the
        literature, not claimed as new;
  \item \textsf{[MacBeth]} --- claimed and/or proved by MacBeth, informally
        (not machine-checked);
  \item \textsf{[MacBeth, Lean-verified: \textit{file}]} --- machine-checked in
        the repository's \texttt{lean/} tree, in the named source file;
  \item \textsf{[Conjecture, MacBeth]} or \textsf{[Open]} --- for conjectures.
\end{itemize}
A tag \textsf{[Cited:\,\ldots]} on a statement that MacBeth has also verified
independently will say so; when in doubt the weaker (cited, not-checked) tag is
used, and a footnote records the uncertainty. The bias is deliberately towards
\emph{under}-claiming. The reading and citation tracker in
Appendix~\ref{app:tracker} records, paper by paper, how thoroughly each source
has been engaged.

A note on the Lean tags. Only two Lean \emph{source} files are committed to the
repository at the time of this draft: \texttt{Basic.lean} (the container
representation theorem) and \texttt{Directed.lean} (the comonad laws from the
directed-container laws). Results checked in those files carry the
\textsf{Lean-verified} tag naming the committed file. Where MacBeth has a Lean
formalisation that is \emph{not} yet committed as source to this repository
(notably the $\DCont\cong\Cof$ and Zappa--Sz\'ep associativity developments),
the statement is tagged \textsf{[Cited]} for the underlying theorem and a
footnote flags that the Lean source is not in this tree --- so the reader is
never told something is machine-checked here when its source cannot be found
here.

\mainmatter

\chapter{Containers and their extension functor}
\label{ch:cont}

\section{Containers}

\begin{Definition}[Container]\label{def:container}
A \emph{container} is a pair $\cont SP$ where $S$ is a set of \emph{shapes} and
$P\colon S\to\Set$ is a family of \emph{positions}, one set $P(s)$ for each shape
$s$. Its \emph{extension} (or \emph{interpretation}) is the functor
$\Ext{\cont SP}\colon\Set\to\Set$ defined on objects by
\[
  \Ext{\cont SP}(X) \;=\; \coprod_{s\in S} X^{P(s)}
    \;=\; \bigl\{ (s,v) \mid s\in S,\ v\colon P(s)\to X \bigr\},
\]
and on a morphism $f\colon X\to Y$ by postcomposition,
$\Ext{\cont SP}(f)(s,v) = (s,\, f\circ v)$.
\end{Definition}

Unpacking: a shape $s$ is a structure with its positions $P(s)$; a labelling
$v\colon P(s)\to X$ fills every position with an element of $X$; the coproduct
over $s$ is the disjoint union ``pick exactly one shape''. The functorial action
relabels --- it changes the data at the positions but never the shape, so the
\emph{shape projection} $\pi(s,v)=s$ is a natural transformation to the constant
functor at $S$. A container thus separates \emph{structure} (the shape) from
\emph{content} (the labels), and its functorial action is exactly \texttt{map}.

\begin{Example}[The container menagerie]\label{ex:menagerie}
\textbf{Maybe}: two shapes, ``nothing'' ($P=\varnothing$, term $X^0=1$) and
``just'' ($P=\{*\}$, term $X$); extension $1+X$.
\textbf{Lists}: shape $n\in\mathbb N$, positions $\{0,\dots,n-1\}$; extension
$\sum_{n} X^n$.
\textbf{Binary trees}: a shape is a planar binary tree, positions are leaves;
the generating function satisfies $T(X)=1+X\cdot T(X)^2$, so the number of shapes
with $n$ positions is the Catalan number $C_n=\frac1{n+1}\binom{2n}{n}$.
\textbf{Matrices}: shape $(m,n)$, positions the $mn$ entries; extension
$\sum_{m,n} X^{m\times n}$. In each case the coefficient of $X^n$ in the
generating function counts \emph{how many} shapes have $n$ positions, while the
container remembers \emph{which} shapes and \emph{where} the positions are.
\end{Example}

A standing intuition (``decategorification''): writing $\sum_s X^{P(s)}$ as if it
were a polynomial in a formal indeterminate $X$ guides the right intuitions about
products, composition and derivatives. But $X$ is a \emph{set}, not an
indeterminate: $X^{P(s)}=\Hom(P(s),X)$ is a set of functions, and the extension
lives in $[\Set,\Set]$, not in a ring of power series.

\section{The representation theorem}

The combinatorial data of containers captures \emph{exactly} the natural
transformations between their extensions.

\begin{Theorem}[Representation / full faithfulness]\label{thm:rep}
The extension functor $\Ext{-}\colon\Cont\to[\Set,\Set]$ is full and faithful:
every natural transformation $\Ext{\cont{S_1}{P_1}}\Rightarrow\Ext{\cont{S_2}{P_2}}$
arises from a unique container morphism, and distinct morphisms induce distinct
natural transformations.
\prov{Cited: Abbott--Altenkirch--Ghani 2005}%
\footnote{Independently formalised by MacBeth in the committed Lean file
\texttt{lean/Containers/Containers/Basic.lean} (round-trip
\texttt{ContainerMorphism}\,$\leftrightarrow$\,\texttt{ExtNatTrans}, no
\texttt{sorry}).}
\end{Theorem}

\begin{proof}[Proof sketch]
By the Yoneda lemma a natural transformation $\alpha$ is determined by its action
on the ``generic element'' $(s,\id_{P_1(s)})\in\Ext{\cont{S_1}{P_1}}(P_1(s))$ of
each shape: $\alpha_{P_1(s)}(s,\id)=(u(s),f_s)$ reads off both a shape map
$u\colon S_1\to S_2$ and a position map $f_s\colon P_2(u(s))\to P_1(s)$, and
naturality forces every other component. Conversely every $(u,f)$ is natural
(Chapter~\ref{ch:contcat}). See \cite{AAG05} for the full argument; the
round-trip is machine-checked in \texttt{Basic.lean}.
\end{proof}

\chapter{The category \texorpdfstring{$\Cont$}{Cont}}
\label{ch:contcat}

\section{Container morphisms: the variance}

\begin{Definition}[Container morphism]\label{def:contmor}
A \emph{container morphism} $(u,f)\colon\cont{S_1}{P_1}\to\cont{S_2}{P_2}$ consists
of
\begin{itemize}[nosep]
  \item a \emph{shape map} $u\colon S_1\to S_2$ (\emph{forward}), and
  \item for each $s\in S_1$, a \emph{position map}
        $f_s\colon P_2(u(s))\to P_1(s)$ (\emph{backward}).
\end{itemize}
It induces the natural transformation $\alpha_X(s,v)=\bigl(u(s),\, v\circ f_s\bigr)$.
\end{Definition}

The variance is not a convention but a necessity. To build a new labelling of the
\emph{output} positions $P_2(u(s))$ out of an old labelling $v\colon P_1(s)\to X$
without inspecting labels, the only move is to precompose:
\[
  P_2(u(s)) \xrightarrow{\ f_s\ } P_1(s) \xrightarrow{\ v\ } X,
\]
which forces $f_s$ to run from output positions to input positions. Each output
position $p'$ asks ``which input position supplies my label?'' and the answer is
$f_s(p')$. Consequently a container morphism can only \emph{rearrange},
\emph{discard} or \emph{duplicate} labels --- never \emph{create} one. (Sorting,
filtering and thresholding are \emph{not} container morphisms: they inspect
labels, and so are not natural in $X$.)

\begin{Definition}[Identity and composition]\label{def:contcomp}
The identity on $\cont SP$ is $(\id_S,\id)$. The composite of
$(u,f)\colon\cont{S_1}{P_1}\to\cont{S_2}{P_2}$ and
$(u',f')\colon\cont{S_2}{P_2}\to\cont{S_3}{P_3}$ has shape map $u'\circ u$ and
position map running backward through both stages,
\[
  (u'\circ u)^\sharp_s \;=\; f_s\circ f'_{u(s)}\colon
   P_3\bigl(u'(u(s))\bigr) \xrightarrow{\ f'_{u(s)}\ } P_2(u(s))
   \xrightarrow{\ f_s\ } P_1(s).
\]
These satisfy the category axioms; the resulting category is $\Cont$.
\end{Definition}

\begin{Proposition}[Naturality and the correspondence]\label{prop:nat}
For every container morphism $(u,f)$, the family $\alpha_X(s,v)=(u(s),v\circ f_s)$
is natural in $X$; and (Theorem~\ref{thm:rep}) every natural transformation
between container extensions is $\alpha$ for a unique $(u,f)$. Hence
$\Ext{-}\colon\Cont\to[\Set,\Set]$ is an isomorphism onto its image.
\prov{Cited: AAG05}
\end{Proposition}

\begin{proof}
Naturality is associativity of composition: for $g\colon X\to Y$, both ways round
the naturality square send $(s,v)$ to $(u(s),\, g\circ v\circ f_s)$, because
$f_s$ does not look at labels. Full faithfulness is Theorem~\ref{thm:rep}.
\end{proof}

\section{Worked example: \texttt{tail}}

The function $\mathrm{tail}\colon\mathrm{List}(X)\to 1+\mathrm{List}(X)$ is a
container morphism $(u,f)$: the shape map sends length $n$ to length $n-1$ (and
$0$ to ``nothing''), and the (backward) position map sends position $i$ of the
output to position $i+1$ of the input. The same $(u,f)$ acts on lists of integers,
strings, or anything --- the labels are shifted one place left, never inspected.
This is the content of naturality: the morphism is parametric in the label type.

\chapter{Algebraic structure on \texorpdfstring{$\Cont$}{Cont}}
\label{ch:algebra}

This chapter is, by design, partly an \emph{outline with cited statements}. The
monoidal and closed structures on containers are most cleanly developed in the
polynomial-functor setting, where $\Cont$ embeds as the full subcategory $\Poly$
of polynomial endofunctors. We record the structures, cite their sources, and
flag carefully which MacBeth has personally verified. We do \emph{not} reprove
them here, and we do not invent proofs.

\begin{Remark}[Containers and polynomial functors]
A container $\cont SP$ presents the polynomial functor
$X\mapsto\sum_{s\in S}X^{P(s)}$; the category $\Poly$ of polynomial functors and
natural transformations is equivalent to $\Cont$ (this is
Theorem~\ref{thm:rep} read as ``$\Poly\simeq\Cont$''). We use the two names
interchangeably and adopt Spivak's catalogue of monoidal structures on $\Poly$
\cite{Spivak21}.
\end{Remark}

\section{Products and coproducts}

\begin{Proposition}[Coproducts in $\Cont$]\label{prop:coprod}
$\Cont$ has all coproducts: the coproduct of $\cont{S_i}{P_i}$ has shape set
$\coprod_i S_i$ and positions inherited shapewise,
$P\bigl((i,s)\bigr)=P_i(s)$. On extensions it is the pointwise coproduct of
functors, $\Ext{\coprod_i C_i}=\coprod_i\Ext{C_i}$.
\prov{Cited: AAG05; Spivak 2021}
\end{Proposition}

\begin{Proposition}[Products in $\Cont$]\label{prop:prod}
$\Cont$ has all products: the product of $\cont{S_1}{P_1}$ and $\cont{S_2}{P_2}$
has shape set $S_1\times S_2$ and positions
$P\bigl((s_1,s_2)\bigr)=P_1(s_1)\sqcup P_2(s_2)$, so that
$\Ext{C_1\times C_2}(X)=\Ext{C_1}(X)\times\Ext{C_2}(X)$.
\prov{Cited: Spivak 2021}%
\footnote{Stated as in \cite{Spivak21}; the coproduct (Prop.~\ref{prop:coprod})
MacBeth has checked directly on extensions, the product MacBeth has \emph{not}
independently re-derived in full and reports it as cited.}
\end{Proposition}

\section{The monoidal menagerie}

Spivak \cite{Spivak21} catalogues several monoidal structures on $\Poly$. We list
the ones relevant to this book; the reader should treat each as
\emph{cited, not independently checked} unless tagged otherwise.

\begin{Definition}[The four products, after Spivak 2021]\label{def:menagerie}
On $\Poly\simeq\Cont$ there are, among others:
\begin{itemize}[nosep]
  \item the \emph{coproduct} $+$ (Prop.~\ref{prop:coprod}), unit the initial
        polynomial $0$;
  \item the \emph{product} $\times$ (Prop.~\ref{prop:prod}), unit the terminal
        polynomial $1$;
  \item the \emph{Dirichlet (parallel) tensor} $\otimes$: on shapes
        $S_1\times S_2$, on positions $P_1(s_1)\times P_2(s_2)$, so that
        $\Ext{C_1\otimes C_2}(X)=\sum_{s_1,s_2}X^{P_1(s_1)\times P_2(s_2)}$; unit
        the identity polynomial $X\mapsto X$ (\,$1\triangleleft 1$\,);
  \item the \emph{composition (substitution) product} $\triangleleft$, with
        $\Ext{C_1\triangleleft C_2}=\Ext{C_1}\circ\Ext{C_2}$; unit the identity
        polynomial. This is the \emph{sequencing} product, and comonoids for it
        are the subject of Chapter~\ref{ch:comonoids}.
\end{itemize}
\prov{Cited: Spivak 2021}
\end{Definition}

\begin{Remark}[Honesty about $\otimes$ vs.\ $\times$]
The symbols are easy to confuse. The \emph{Dirichlet} tensor $\otimes$ multiplies
\emph{position sets} ($P_1\times P_2$); the categorical \emph{product} $\times$
takes a \emph{disjoint union} of position sets ($P_1\sqcup P_2$). They have the
same shape set $S_1\times S_2$ but different position data and different
extensions. Both are due to Spivak's catalogue \cite{Spivak21}; MacBeth has
verified the shape/position formulas combinatorially on small cases but defers
the unit and coherence laws to \cite{Spivak21}.
\end{Remark}

\begin{Definition}[Closed structures]\label{def:closed}
$(\Poly,\otimes)$ and $(\Poly,\times)$ are closed: each has an internal hom
(``$[\,-,-\,]$'') right adjoint to the corresponding tensor. The composition
product $\triangleleft$ is \emph{not} symmetric and its closure is subtler (the
relevant right adjoints exist only in restricted forms).
\prov{Cited: Spivak 2021}%
\footnote{The internal homs are stated as in \cite{Spivak21}; MacBeth has
\emph{not} independently verified the adjunctions and reports them as cited.}
\end{Definition}

The free-monad and cofree-comonad constructions associated with $\triangleleft$
are deferred to Chapter~\ref{ch:phase2}, where they belong to the functors-over-$\Cont$
programme rather than to the structure of $\Cont$ itself.

\chapter{Comonoids in \texorpdfstring{$\Cont$}{Cont}: directed containers and small categories}
\label{ch:comonoids}\label{ch:phase2dummy}

This is the heart of the book. We show that comonoids for the composition product
$\triangleleft$ are \emph{directed containers}, that directed containers are
\emph{small categories}, and --- the refinement that the rest of the grant turns
on --- that the right \emph{morphisms} are \emph{cofunctors}, not functors.

\section{Directed containers}

\begin{Definition}[Directed container]\label{def:dcont}
A \emph{directed container} is a container $\cont SP$ together with
\begin{itemize}[nosep]
  \item a \emph{root} $\mathsf o(s)\in P(s)$ for each $s$;
  \item a \emph{sub-shape} operation $s\downarrow p\in S$ for $p\in P(s)$;
  \item a \emph{shift} $p\oplus q\in P(s)$ for $p\in P(s)$, $q\in P(s\downarrow p)$;
\end{itemize}
subject to the five laws (for all $s$, $p\in P(s)$, $q\in P(s\downarrow p)$,
$q'\in P((s\downarrow p)\downarrow q)$):
\begin{align*}
  &\text{(D1)} & s\downarrow\mathsf o(s)&=s;
  & &\text{(D2)} & \mathsf o(s)\oplus q&=q;
  & &\text{(D3)} & p\oplus\mathsf o(s\downarrow p)&=p;\\
  &\text{(D4)} & s\downarrow(p\oplus q)&=(s\downarrow p)\downarrow q;
  & &\text{(D5)} & (p\oplus q)\oplus q'&=p\oplus(q\oplus q').
\end{align*}
\end{Definition}

\noindent D2, D3, D5 make each $P(s)$ a monoid-like structure with unit
$\mathsf o(s)$; D1, D4 track how sub-shapes nest. (The well-typedness of D5 needs
D4 --- exactly the dependency that recurs as a dependent-cast obstacle in the Lean
formalisation.)

\section{Directed containers are comonads}

\begin{Definition}[The comonad of a directed container]\label{def:dcomonad}
On $D:=\Ext{\cont SP}$ define
\[
  \varepsilon(s,v)=v(\mathsf o(s)),\qquad
  \delta(s,v)=\bigl(s,\ \lambda p.\,(s\downarrow p,\ \lambda q.\, v(p\oplus q))\bigr).
\]
\end{Definition}

\begin{Theorem}[Comonad laws $\Longleftrightarrow$ D1--D5]\label{thm:comonad}
$(D,\varepsilon,\delta)$ is a comonad, and each comonad law is exactly one
directed-container law: left counit is D1 (shape) and D2 (value); right counit is
D3; coassociativity is D4 (shapes) and D5 (deepest values). Conversely a comonad
of this form forces D1--D5.
\prov{Cited: Ahman--Chapman--Uustalu 2014; MacBeth, Lean-verified:
Directed.lean}
\end{Theorem}

\begin{proof}[Proof sketch]
Direct calculation. For left counit, $\varepsilon(\delta(s,v))$ reduces to
$(s\downarrow\mathsf o(s),\,\lambda q.\,v(\mathsf o(s)\oplus q))$, which is $(s,v)$
by D1 then D2. Right counit is D3; coassociativity unfolds both composites to a
common triple nest, equal by D4 (inner shapes) and D5 (deepest values). The full
computation, and the converse, are checked with zero \texttt{sorry} in the
committed file \texttt{lean/Containers/Containers/Directed.lean}
(\texttt{left\_counit}, \texttt{right\_counit}, \texttt{coassoc}); the
``container is a comonad'' direction is Ahman--Chapman--Uustalu \cite{ACU14}.
\end{proof}

\section{Directed containers are small categories}

\begin{Theorem}[Object dictionary]\label{thm:objdict}
Directed-container structures on $\cont SP$ are in bijection with small categories
$\mathcal C$ with object set $S$ and, for each $s$,
$P(s)=\coprod_{s'}\mathcal C(s,s')$ the morphisms \emph{out of} $s$, under
\[
  \ob\mathcal C=S,\quad
  \mathcal C(s,s')=\{p\in P(s):s\downarrow p=s'\},\quad
  \id_s=\mathsf o(s),\quad
  q\circ p=p\oplus q .
\]
In particular $s\downarrow p=\cod(p)$, and every small category arises from a
unique directed container.
\prov{Cited: Ahman--Uustalu 2016}
\end{Theorem}

Combining Theorems~\ref{thm:comonad} and \ref{thm:objdict} with the polynomial
incarnation gives the object-level chain promised in the Preface: a directed
container is a comonoid in $(\Poly,\triangleleft)$, hence
\[
  \textbf{Containers}\ \supseteq\ \textbf{Directed Containers}
  \ \cong\ \textbf{Polynomial Comonoids}\ \cong\ \textbf{Small Categories}.
\]

\section{The morphism refinement: \texorpdfstring{$\DCont\cong\Cof$}{DCont = Cof}}

A dictionary on objects invites a dictionary on morphisms. The naive guess --- a
morphism of directed containers is a \emph{functor} --- is \emph{false on
variance}. A container morphism has a backward position map
$f^\sharp_s\colon P'(fs)\to P(s)$, whereas a functor pushes morphisms forward. The
maps native to directed containers are \emph{cofunctors}.

\begin{Definition}[Cofunctor / retrofunctor]\label{def:cofunctor}
A \emph{cofunctor} $\varphi\colon\mathcal C\co\mathcal D$ is an object map
$\varphi_0\colon\ob\mathcal C\to\ob\mathcal D$ together with a \emph{lift}: for each
object $c$ and each $h\colon\varphi_0 c\to d$ in $\mathcal D$, a morphism
$\varphi^\uparrow(c,h)\colon c\to c'$ in $\mathcal C$, subject to
\begin{align*}
  \text{(C0)}\ & \varphi_0(\cod\varphi^\uparrow(c,h))=\cod(h);\quad
  \text{(C1)}\ \varphi^\uparrow(c,\id_{\varphi_0 c})=\id_c;\\
  \text{(C2)}\ & \varphi^\uparrow(c,h'\circ h)=\varphi^\uparrow(c',h')\circ
        \varphi^\uparrow(c,h),\quad c'=\cod\varphi^\uparrow(c,h).
\end{align*}
Small categories and cofunctors form a category $\Cof$.
\prov{Cited: Aguiar 1997; Clarke 2020; Spivak 2021 (Def.\ 3.21)}
\end{Definition}

\begin{Remark}[$\Cof=\Cathash$]
$\Cof$ is exactly $\Cathash=\mathrm{Comod}(\Poly)$, the category of comonoids in
$(\Poly,\triangleleft)$ and their morphisms: a cofunctor \emph{is} a comonoid
homomorphism \cite{Spivak21,SS2405}. This is why the object dictionary
(Theorem~\ref{thm:objdict}) cannot extend to $\Cat$: functors are carried not by
comonoid morphisms but by bicomodules \cite{SS23}.
\end{Remark}

\begin{Definition}[Directed-container morphism]\label{def:dcm}
A \emph{morphism of directed containers} $D\to D'$ is a container morphism
$(f,f^\sharp)$ with $f\colon S\to S'$ and $f^\sharp_s\colon P'(fs)\to P(s)$,
satisfying, for all $s$, $h\in P'(fs)$, $k\in P'((fs)\downarrow' h)$,
\begin{align*}
  \text{(M0)}\ & f(s\downarrow f^\sharp_s(h))=(fs)\downarrow' h;\quad
  \text{(M1)}\ f^\sharp_s(\mathsf o'(fs))=\mathsf o(s);\\
  \text{(M2)}\ & f^\sharp_s(h\oplus' k)=
        f^\sharp_s(h)\ \oplus\ f^\sharp_{\,s\downarrow f^\sharp_s(h)}(k).
\end{align*}
\end{Definition}

\begin{Theorem}[The morphism dictionary]\label{thm:dcontcof}
The assignment $(f,f^\sharp)\mapsto(\varphi_0:=f,\ \varphi^\uparrow(s,h):=f^\sharp_s(h))$
is a bijection on every hom-set carrying (M0)--(M2) to (C0)--(C2) term for term,
and it respects identities and composition. It is moreover a bijection on objects
(Theorem~\ref{thm:objdict}). Hence there is an isomorphism of categories
\[
  \DCont\ \cong\ \Cof .
\]
\prov{Cited: Ahman--Uustalu 2016 \S3.2--3; Shapiro--Spivak 2023, Cor.\,5.12}%
\footnote{MacBeth has formalised this isomorphism in Lean (file
\texttt{Cofunctor.lean}, built and reported zero-\texttt{sorry} in
\texttt{papers/dcont-cof.tex}), but that source is \emph{not yet committed} to
this repository's \texttt{lean/} tree; so the statement is tagged \textsf{[Cited]}
here, not \textsf{Lean-verified}, pending the source landing in the repo. The
\emph{identification} of Ahman--Uustalu's ``relative split pre-opcleavages'' with
Spivak's cofunctors, and the bridge to delta lenses below, are MacBeth's
connective contribution; see \texttt{papers/dcont-cof.tex}.}
\end{Theorem}

\begin{proof}[Proof sketch]
Under the object dictionary a position $h\in P'(fs)$ is a morphism out of $fs$ and
$f^\sharp_s(h)$ a morphism out of $s$, so $(f,f^\sharp)$ and
$(\varphi_0,\varphi^\uparrow)$ are literally the same data. Reading $\oplus$ as
reverse composition, (M0)$\Leftrightarrow$(C0) (codomain), (M1)$\Leftrightarrow$(C1)
(unit), (M2)$\Leftrightarrow$(C2) (composition). Identities and composites match,
and objects are in bijection by Theorem~\ref{thm:objdict}; a bijective-on-objects,
fully faithful functor is an isomorphism of categories. The full term-by-term
match is in \texttt{papers/dcont-cof.tex} (Lemma~2.5--Theorem~2.7 there).
\end{proof}

\begin{Proposition}[Functors are the opposite variance]\label{prop:cov}
The \emph{covariant} position maps $\phi_s\colon P(s)\to P'(fs)$ subject to the
mirror conditions (N0)--(N2) of (M0)--(M2) are in bijection with \emph{functors}
$\mathcal C(D)\to\mathcal C(D')$. Thus a directed structure carries two notions of
map of opposite variance: covariant gives functors, contravariant gives
cofunctors, and only the latter is a morphism of the underlying container.
\prov{MacBeth}%
\footnote{Verified combinatorially over $36$ ordered pairs of small categories
drawn from $\{\mathbf1,\mathbf2,\mathbf3,\mathbb Z/2,\mathbb Z/3,M\}$;
\#DCont-morphisms $=$ \#cofunctors in all $36$ cases, \#op-variance-morphisms $=$
\#functors in all $36$, and \#functors $\neq$ \#cofunctors in $17$ of the $36$.
See \texttt{papers/dcont-cof.tex} \S5. The underlying fact that functors are
carried by bicomodules, not comonoid morphisms, is \cite{SS23}.}
\end{Proposition}

\begin{Corollary}[Cofunctors are the Put of a delta lens]\label{cor:lens}
A \emph{delta lens} $A\rightleftarrows B$ is a \emph{Get} functor together with a
\emph{Put} cofunctor sharing one object map, related by the section condition
(PutGet) $f\cdot\varphi(a,u)=u$. Under $\DCont\cong\Cof$ the Put cofunctor is
exactly a morphism of directed containers; so a directed-container morphism is the
update-propagating half of a bidirectional transformation.
\prov{Cited: Clarke 2022 (arXiv:2108.00390), Def.\ 4 / Thm.\ 9}%
\footnote{The section condition is (PutGet) $f\cdot\varphi(a,u)=u$, \emph{not}
``$f\circ\varphi=\id$''; MacBeth flags this because the loose form recurs in the
folklore.}
\end{Corollary}

\chapter{Composing systems: Zappa--Sz\'ep and distributive laws}
\label{ch:zs}

Chapter~\ref{ch:comonoids} settled what a single directed container is and what
maps it has. This chapter asks how \emph{two} of them compose into one --- the
question with the economic consequences (mis-composed supply chains lose liability
traceability; a wrong contract composite is an exploit). We package this as a
\emph{Zappa--Sz\'ep} (two-sided semidirect / matched-pair) product of small
categories, the categorical generalisation of the group/monoid Zappa--Sz\'ep
product.

\section{The pairwise criterion}

For objects $a,b$ of a small category, $\Hom(a,b)$ is canonically a right
$\End(a)$-set (precomposition) and a left $\End(b)$-set (postcomposition). A
Zappa--Sz\'ep decomposition $K=\mathcal C\bowtie\mathcal D$ of $K$ over a wide
subcategory $\mathcal D$ asks for a \emph{unique factorisation} of each morphism
through a transversal and an endomorphism.

\begin{Theorem}[Pairwise Zappa--Sz\'ep criterion]\label{thm:zscriterion}
Let $\mathcal D$ be a wide subcategory of $K$. Then $K=\mathcal C\bowtie\mathcal D$
for a complementary wide subcategory $\mathcal C$ if and only if
\begin{itemize}[nosep]
  \item \textup{(L)} each hom-presheaf $\Hom_K(-,b)$ is a \emph{free} right
        $\mathcal D$-set (equivalently: the source-oriented factorisation
        $f=t\circ e$ is unique for every $f$), and
  \item \textup{(G)} the free bases close into a wide subcategory $\mathcal C$.
\end{itemize}
\prov{MacBeth}%
\footnote{This is the corrected form of the SEED conjecture ``every small category
is a Zappa--Sz\'ep product of a poset skeleton and its vertex monoids,'' which was
\emph{refuted} computationally (smallest counterexample: two objects, one arrow
$s\colon a\to b$, $\End(a)=\End(b)=\mathbb Z/2$ fixing $s$, so
$2\nmid 1=|\Hom(a,b)|$). The freeness half (L) is written up self-contained in
\texttt{papers/pairwise-zappa-szep.tex} (a built PDF is in the repo; the
\texttt{.tex} source is not yet committed to this tree). A brute-force enumeration
of $\approx 26$ small categories matches the criterion exactly.}
\end{Theorem}

\begin{Theorem}[Cocycle axioms $\Longleftrightarrow$ associativity]\label{thm:zsassoc}
The Zappa--Sz\'ep cocycle axioms ZS1--ZS4 (relating the two actions of a matched
pair) hold if and only if the product on $\mathcal C\bowtie\mathcal D$ is
associative.
\prov{MacBeth}%
\footnote{MacBeth has a Lean formalisation of the monoid case (ZS1--ZS4
$\Longleftrightarrow$ associativity of the Zappa--Sz\'ep product) in
\texttt{ZappaSzep.lean}, built and reported zero-\texttt{sorry}; that source is
\emph{not yet committed} to this repository's \texttt{lean/} tree, so the tag is
\textsf{[MacBeth]} rather than \textsf{Lean-verified} here.}
\end{Theorem}

\section{(G) as a strict factorisation system, and its $H^2$ obstruction}

\begin{Proposition}[(G) is a strict factorisation system]\label{prop:sfs}
Condition (G) --- that the free bases close into a wide subcategory --- is exactly
the statement that the decomposition extends to a \emph{strict factorisation
system} on $K$. Distributive laws of categories correspond to strict factorisation
systems.
\prov{Cited: Rosebrugh--Wood (distributive laws $\Longleftrightarrow$ SFS)}
\end{Proposition}

\begin{Theorem}[The (G) obstruction is an $H^2$ class]\label{thm:h2}
Under a trivial-action hypothesis, (G) is classified by a second cohomology class:
$\textup{(G)}\Longleftrightarrow[\omega]=0$, where $[\omega]$ lives in the
\emph{Baues--Wirsching} cohomology $H^2(\mathrm{Sk}_{\mathcal C};\mathcal D)$ of the
$\mathcal C$-skeleton with a source-only (presheaf) natural system of coefficients.
This is the split case of the Baues--Wirsching linear-extension classification.
\prov{MacBeth (identification), resting on Cited: Baues--Wirsching 1985}%
\footnote{The cohomological \emph{identification} --- that the obstruction to (G)
is a Baues--Wirsching $H^2$ class and that ``(G)$\Leftrightarrow[\omega]=0$'' is the
split case of the BW linear-extension theorem --- is MacBeth's; the underlying
classification theorem is Baues--Wirsching, \emph{Cohomology of small categories},
JPAA \textbf{38} (1985) 187--211 (modern account: arXiv:2508.00727). The rigid
twist of the counterexample above is the $\mathbb Z/2$ generator; the nonabelian
boundary is the Kac/Rosebrugh--Wood law.}
\end{Theorem}

\section{The classical anchor}

\begin{Proposition}[Classical Zappa--Sz\'ep recovered]\label{prop:monoidanchor}
Specialising Theorem~\ref{thm:zscriterion} to a one-object category (a monoid)
recovers the classical Zappa--Sz\'ep product of monoids, and (L) becomes the
familiar condition that the factorisation map is a bijection.
\prov{MacBeth}%
\footnote{MacBeth flags an earlier \emph{over}claim for correction: the
``Ahman--Uustalu update monad as a degenerate Zappa--Sz\'ep product'' reading is
\emph{not} a special case of the two-category matched pair --- it is a monoid acting
on a set, not two composable monoids --- and should not be presented as such.}
\end{Proposition}

\chapter{Functors over the category of containers (Phase 2 --- outline)}
\label{ch:phase2}

This chapter is \emph{mostly} a stub sketching Neil's proposed Phase 2 programme; the one
exception is \S\ref{sec:freecofree}, which contains worked, small-case-verified results on the
free monad and cofree comonad. Where Chapters~\ref{ch:cont}--\ref{ch:zs} studied containers
and the maps \emph{between} them, Phase 2 studies \emph{endofunctors of $\Cont$}
--- operations $\Cont\to\Cont$ (equivalently $\Poly\to\Poly$) and the structures
they carry.

\paragraph{What a functor $\Cont\to\Cont$ is.} It assigns to each container a
container, functorially in container morphisms. The programme is to identify the
natural such operations and the (co)monad structures they support.

\paragraph{Candidate examples (to be developed, not yet results).}
\begin{itemize}[nosep]
  \item \emph{The derivative / one-hole-context functor} $\partial$: sends
        $\cont SP$ to the container of ``a shape with one position singled out''.
        For lists, $\partial(\sum_n X^n)$ is the zipper. The combinatorial
        derivative is classical for species; its container-level naturality and
        its interaction with $\triangleleft$ are the targets here.
  \item \emph{The free monad functor} on $(\Poly,\triangleleft)$ and the dual
        \emph{cofree comonad functor}: the syntax/behaviour pair. These two are
        developed concretely in \S\ref{sec:freecofree} below --- their shape/position
        data and (co)monad laws are worked out and small-case verified, so unlike the
        derivative bullet they are \emph{results}, not stubs.
\end{itemize}

\section{The free monad and cofree comonad of a container}
\label{sec:freecofree}

We work where W- and M-types exist (e.g.\ $\Set$). For a container $C=\cont SP$ both
the free monad and the cofree comonad of $\Ext C$ are again containers; we give their
shape/position data explicitly and reduce the (co)monad laws to the monoid laws of
\emph{tree grafting} and \emph{path concatenation}. Throughout, $I=\cont 11$ is the
$\triangleleft$-unit.

\begin{Definition}[Free monad container]\label{def:freemonad}
$C^{*}=\cont{S^{*}}{P^{*}}$ where
$S^{*}=\mu Y.\,\bigl(1+\textstyle\sum_{s:S}(P(s)\to Y)\bigr)$ is the set of well-founded
$C$-trees with a distinguished \emph{variable-leaf} constructor (the $1$), and $P^{*}(w)$
is the set of variable leaves of $w$:
$P^{*}(\bullet)=1$, $P^{*}(\mathsf{node}(s,c))=\sum_{p:P(s)}P^{*}(c(p))$.
\end{Definition}

\begin{Definition}[Cofree comonad container]\label{def:cofreecomonad}
$C^{\infty}=\cont{S^{\infty}}{P^{\infty}}$ where
$S^{\infty}=\nu Y.\,\textstyle\sum_{s:S}(P(s)\to Y)$ is the M-type of (possibly infinite)
$C$-trees, and $P^{\infty}(w)$ is the set of \emph{nodes} of $w$, i.e.\ finite paths from the
root: $P^{\infty}(\langle s,c\rangle)=1+\sum_{p:P(s)}P^{\infty}(c(p))$ (the empty path
$\bullet$ plus the paths into each child). Each path is finite even when $w$ is infinite.
\end{Definition}

\begin{Proposition}[Extensions]\label{prop:freecofree-ext}
$\Ext{C^{*}}A\cong\mu X.\,(A+\Ext C X)$ is the carrier of the free monad on $\Ext C$, and
$\Ext{C^{\infty}}A\cong\nu X.\,(A\times\Ext C X)$ is the carrier of the cofree comonad on
$\Ext C$ --- the possibly-infinite $C$-trees with \emph{every node} labelled in $A$.
\prov{Cited: Abbott--Altenkirch--Ghani 2005 (W/M-types); Spivak 2021 (tree polynomial).
Container-language derivation re-done here; small cases below.}
\end{Proposition}

The structure maps are container morphisms.
\emph{Unit} $\eta:I\to C^{*}$ sends $\ast$ to the variable leaf $\bullet$.
\emph{Multiplication} $\mu:C^{*}\triangleleft C^{*}\to C^{*}$ grafts the tree $g(p)$ into each
variable leaf $p$ of $w$; on positions it is the canonical iso
$P^{*}(\mathsf{graft}(w,g))\cong\sum_{p:P^{*}(w)}P^{*}(g(p))$, i.e.\ the identity.
Dually, \emph{counit} $\varepsilon:C^{\infty}\to I$ has $f_{\varepsilon}(w)(\ast)=\bullet$
(the root), and \emph{comultiplication} $\delta:C^{\infty}\to C^{\infty}\triangleleft
C^{\infty}$ has $u_{\delta}(w)=(w,\lambda n.\,w/n)$ (label each node $n$ by the subtree $w/n$)
and $f_{\delta}(w)(n,m)=n\cdot m$ (path concatenation).

\begin{Theorem}[(Co)monad laws via grafting and concatenation]\label{thm:freecofree-laws}
$(C^{*},\eta,\mu)$ is a monoid and $(C^{\infty},\varepsilon,\delta)$ a comonoid in
$(\Cont,\triangleleft)$. The monad laws reduce to the monoid laws of grafting (the variable
leaf is neutral; substitution is associative); the comonad laws reduce to those of path
concatenation (the root is the empty path; concatenation is associative).
\prov{MacBeth, 2026-06-11: laws checked by exhaustive enumeration of small labelled trees for
$C=1+X$ (Maybe) and $C=1+X^{2}$ (binary); scripts in \texttt{projects/scratch/}.}
\end{Theorem}

\begin{Theorem}[$C^{\infty}$ is a directed container]\label{thm:cofree-dircont}
$C^{\infty}$ is a directed container with root $\mathsf o(w)=\bullet$, sub-shape
$w\downarrow n=w/n$ (the subtree at node $n$), and shift $n\oplus m=n\cdot m$. Laws
\textup{D1--D5} (Def.~\ref{def:dcont}) hold, the only non-formal one being
$w/(n\cdot m)=(w/n)/m$ behind \textup{D4}. By Theorem~\ref{thm:comonad} this directed
container \emph{is} the comonad above: the cofree comonad of a container is the cofree
\emph{directed} container on it. Its category (Theorem~\ref{thm:objdict}) is the
\emph{subtree category}: objects are trees, a morphism $w\to w'$ is a node $n$ of $w$ with
$w/n=w'$, identities are roots, composition is concatenation.
\prov{Cited in substance: Ahman--Chapman--Uustalu (polynomial comonad $=$ category); Spivak,
and Spivak--Niu \emph{Polynomial Functors} (cofree comonoid $=$ tree comonoid). MacBeth:
explicit $\mathsf o/\!\downarrow\!/\oplus$ data in the D1--D5 form and the verification.}
\end{Theorem}

\begin{Example}[Maybe and binary]\label{ex:freecofree}
For $C=1+X$ (shapes $\mathsf{nil}$ with $P=\varnothing$, $\mathsf{cons}$ with $P=1$):
$S^{\infty}=\nu Y.\,1+Y=\bar{\mathbb N}$ (conaturals) and $\Ext{C^{\infty}}A$ is the
nonempty colists of $A$ --- the classical $\mathsf{head}/\mathsf{tails}$ comonad; while
$\Ext{C^{*}}A$ is the free monad on $\mathsf{Maybe}$ (lists terminated by a variable or
$\mathsf{nil}$). For the binary container $C=1+X^{2}$, $C^{*}$ is finite binary trees with
variable leaves and $C^{\infty}$ is possibly-infinite $A$-labelled binary trees, $\delta$
relabelling each node by its subtree.
\end{Example}

\paragraph{Why it matters for the grant.} Free monads and cofree comonads are the
syntax/behaviour pair for compositional systems: a system's available moves
(free monad) and its observable unfolding (cofree comonad). Verified composition
at this level would let one certify, e.g., that an agent-orchestration combinator
preserves the lens (Put/Get) structure of its parts. This is flagged as the
\emph{next direction}; nothing in this chapter is claimed as proved.

\appendix
\chapter{Reading \& citation tracker}
\label{app:tracker}

The status column records how thoroughly MacBeth has engaged each source: ``read
in full'' (worked through, definitions and key proofs internalised), ``skimmed''
(read for the relevant statement), or ``to-read''.

\begin{center}
\renewcommand{\arraystretch}{1.25}
\begin{tabular}{@{}p{5.4cm}p{3.1cm}p{2.4cm}p{2.6cm}@{}}
\toprule
\textbf{Paper} & \textbf{arXiv / ref} & \textbf{Status} & \textbf{Used in} \\
\midrule
Abbott--Altenkirch--Ghani, \emph{Containers: constructing strictly positive types} (2005)
  & TCS 342 & read in full & Ch.\,1--2 \\
Ahman--Chapman--Uustalu, \emph{When is a container a comonad?} (2014)
  & LMCS 10(3) & read in full & Ch.\,4 \\
Ahman--Uustalu, \emph{Directed containers as categories} (2016)
  & arXiv:1604.01187 & read in full & Ch.\,4 \\
Spivak, \emph{Functorial aggregation / Polynomial functors} (2021)
  & arXiv:2111.10968 & read in full & Ch.\,3--4 \\
Shapiro--Spivak, \emph{Structures on categories of polynomials} (2023)
  & arXiv:2305.00167 & read in full & Ch.\,4 \\
Shapiro--Spivak, \emph{A polynomial construction of nerves} (2024)
  & arXiv:2405.13157 & skimmed & Ch.\,4 \\
Clarke, \emph{Internal lenses as functors and cofunctors} (2020)
  & arXiv:2009.06835 & read in full & Ch.\,4 \\
Clarke, \emph{Delta lenses as coalgebras for a comonad} (2022)
  & arXiv:2108.00390 & read in full & Ch.\,4 \\
Rosebrugh--Wood, \emph{Distributive laws and factorization} (JPAA 175, 2002)
  & JPAA 175 & read in full & Ch.\,5 \\
Baues--Wirsching, \emph{Cohomology of small categories} (1985)
  & JPAA 38, 187--211 & read in full & Ch.\,5 \\
\emph{(modern BW account)} (2025)
  & arXiv:2508.00727 & skimmed & Ch.\,5 \\
Aguiar, \emph{Internal categories and quantum groups} (1997)
  & PhD, Cornell & skimmed & Ch.\,4 \\
\bottomrule
\end{tabular}
\end{center}

\begin{thebibliography}{99}

\bibitem{AAG05} M.~Abbott, T.~Altenkirch, N.~Ghani.
\emph{Containers: constructing strictly positive types}.
Theoretical Computer Science \textbf{342} (2005) 3--27.

\bibitem{ACU14} D.~Ahman, J.~Chapman, T.~Uustalu.
\emph{When is a container a comonad?}
Logical Methods in Computer Science \textbf{10}(3), 2014.

\bibitem{AU16} D.~Ahman, T.~Uustalu.
\emph{Directed containers as categories}.
Proc.\ MSFP 2016, EPTCS \textbf{207}, 89--98.
\href{https://arxiv.org/abs/1604.01187}{arXiv:1604.01187}.

\bibitem{Spivak21} D.~I.~Spivak.
\emph{Functorial aggregation} (polynomial functors).
\href{https://arxiv.org/abs/2111.10968}{arXiv:2111.10968}, 2021.

\bibitem{SS23} B.~Shapiro, D.~I.~Spivak.
\emph{Structures on categories of polynomials}.
\href{https://arxiv.org/abs/2305.00167}{arXiv:2305.00167}, 2023.

\bibitem{SS2405} B.~Shapiro, D.~I.~Spivak.
\emph{A polynomial construction of nerves for higher categories}.
\href{https://arxiv.org/abs/2405.13157}{arXiv:2405.13157}, 2024.

\bibitem{Clarke20} B.~Clarke.
\emph{Internal lenses as functors and cofunctors}.
Proc.\ ACT 2019, EPTCS \textbf{323}, 183--195.
\href{https://arxiv.org/abs/2009.06835}{arXiv:2009.06835}.

\bibitem{Clarke22} B.~Clarke.
\emph{Delta lenses as coalgebras for a comonad}.
\href{https://arxiv.org/abs/2108.00390}{arXiv:2108.00390}, 2022.

\bibitem{Aguiar} M.~Aguiar.
\emph{Internal categories and quantum groups}.
PhD thesis, Cornell University, 1997.

\bibitem{RW} R.~Rosebrugh, R.~J.~Wood.
\emph{Distributive laws and factorization}.
Journal of Pure and Applied Algebra \textbf{175} (2002) 327--353.

\bibitem{BW85} H.-J.~Baues, G.~Wirsching.
\emph{Cohomology of small categories}.
Journal of Pure and Applied Algebra \textbf{38} (1985) 187--211.

\bibitem{BWmodern} \emph{(Modern account of Baues--Wirsching cohomology)}.
\href{https://arxiv.org/abs/2508.00727}{arXiv:2508.00727}, 2025.

\end{thebibliography}

\end{document}
